\documentclass[11pt,a4paper]{article}
\usepackage[utf8]{inputenc}
\usepackage[T1]{fontenc}
\usepackage{estilo_capstone}
\setlength{\emergencystretch}{3em}
\sloppy
\titulo{Manual técnico y de despliegue — Evento Cultural}
\subtitulo{Fase 3: reproducción y límites operativos}
\autor{Roberto Riffo · José Miguel Moncada · David Herrera}
\fecha{31 de octubre de 2026}
\begin{document}
\maketitlecapstone
\section{Alcance}
Guía para el equipo técnico. Fuente principal: copia local
\nolinkurl{EventoCultural_Integrado/EventoCultural_Integrado}; referencia remota:
\url{https://github.com/sunsitox/Proyecto_Capston_Evento_Cultural/tree/deploy/vercel-preview}.
No se certifica igualdad con la rama remota ni se despliega al generar este documento.
La SPA usa Supabase directamente para catálogo y sesión; las operaciones privilegiadas
se ejecutan en funciones Edge TypeScript. La carpeta Python \code{backend/} es legado,
no un servicio necesario para iniciar la web actual.

\section{Reproducir el frontend}
Desde el directorio que contiene \code{package.json} y \code{package-lock.json},
utilizar Node/npm compatibles con Angular 22 y las dependencias del lockfile.
Registrar las versiones de Node y npm del entorno de entrega; este documento no
inventa un resultado en una máquina limpia.
\begin{verbatim}
node --version
npm --version
npm ci
npm start
# Controles de calidad, según el cambio:
npm test -- --watch=false
npm run build
npm run lint
\end{verbatim}
\code{npm start} ejecuta \code{ng serve}; por defecto se consulta
\url{http://localhost:4200}. La salida del build es \code{www}, según
\code{angular.json}. Un build aprobado no prueba correo, RLS ni OAuth.
El catálogo de demostración de respaldo debe distinguirse del catálogo remoto.

\section{Dependencias de Supabase}
\begin{itemize}
\item Configurar URL y clave pública en los entornos Angular existentes. Solo
configuración pública puede terminar en el bundle del navegador.
\item Verificar políticas RLS, perfiles y permisos por usuario. El rol administrativo
se consulta en \code{public.profiles.role}, no se concede mediante metadata editable.
\item Aplicar de manera controlada las migraciones necesarias, incluida
\code{20261024000100\_password\_reset\_tokens.sql}. Contrastar el esquema real con
las migraciones antes de modificarlas; documentar respaldo y resultado.
\item Preparar Storage \code{avatars} y sus políticas de propiedad para imágenes.
No conceder escritura global para resolver un error de permisos.
\item Desplegar las funciones presentes en \code{supabase/functions/} mediante el
procedimiento del proyecto y verificar su configuración de autenticación:
\code{registrar-cuenta}, \code{admin-account}, \code{publicar-evento},
\code{admin-eventos}, \code{reset-password-request},
\code{reset-password-confirm}, \code{recomendaciones-semanal}.
\end{itemize}

Los secretos se configuran en el entorno del servidor, nunca en Angular ni Git:
Supabase de servicio, \code{RESEND\_API\_KEY}, firma de recuperación y origen permitido
según los nombres leídos en cada función. Usar las instrucciones del código y no
copiar valores privados al informe. El registro principal crea una cuenta confirmada,
abre sesión y solicita bienvenida por Resend. Si la función no está disponible, el
respaldo legacy puede exigir verificación por correo. No declarar el flujo nuevo
operativo hasta probar la función desplegada.

La recuperación actual genera enlace propio a
\code{/restablecer-contrasena?token=...}; depende de solicitud, firma, tabla de tokens,
correo y confirmación. Validar caducidad y reutilización sin conservar tokens en capturas.
Los antiguos pasos de SMTP/GoTrue no describen por sí solos este flujo.

\section{Vercel}
La configuración ya incluye \code{vercel.json} con reescritura a
\code{/index.html}; no es un archivo pendiente de crear.
Configurar el directorio raíz del proyecto que contiene \code{package.json},
compilación \code{npm run build} y salida \code{www}; conservar el lockfile.
Registrar URL, rama y SHA del despliegue autorizado. El antecedente QA es
\url{https://eventocultural.vercel.app}, ciclo del 11-10-2026, no una comprobación actual.
Verificar navegación directa a ficha, callback y rutas internas; inicio/cierre de sesión;
y respuestas de funciones. La aplicación no requiere publicar la API Python antigua.

\section{Google OAuth y administración}
Google utiliza Supabase Auth y callback \code{/auth/callback}; comprobar proveedor,
Redirect URLs de desarrollo/producción y regreso nativo según \code{OAUTH\_SETUP.md}.
La presencia del botón no demuestra que el proveedor esté configurado.
El panel se abre en \code{/v00/admin}. Solo cuentas de proveedor email con rol
\code{admin} pueden entrar. \code{admin-eventos} también valida identidad y rol:
no resolver un rechazo eliminando comprobaciones ni creando contraseñas fijas.
El cierre por inactividad ocurre en cliente tras 10 minutos, revisados cada 15 segundos;
una recarga inicia un nuevo seguimiento. Probar ese límite expresamente.

\section{Android}
Existe la plataforma \code{android/} y un APK debug local. No volver a agregarla.
\begin{verbatim}
npm run build
npx cap sync android
npx cap open android
\end{verbatim}
Compilar en Android Studio con SDK/JDK compatibles con el proyecto Gradle existente.
Identidad según \code{capacitor.config.ts}: \code{cl.eventocultural.app}, salida
\code{www}. Registrar hash y fecha del APK entregado; la mera existencia del APK no
garantiza correspondencia con el código actual. Probar instalación, permisos,
ubicación denegada/GPS, OAuth y persistencia en teléfono. No versionar firma privada.

\section{Scraper y automatización}
El scraper mantiene ejecución independiente en su repositorio, con dependencias y
secretos propios. Su comando de pruebas documentado es
\code{python -m unittest discover -s tests -v}. El ciclo CI histórico demuestra
extracción/sincronización y conservación de artefactos; no acredita una ejecución hoy.
Verificar el cron real en el workflow: no prometer ejecución puntual a las 23:37
America/Santiago, especialmente ante cambios horarios y retrasos de GitHub Actions.
NVIDIA y la clave de servicio permanecen fuera del cliente.
Docker no forma parte de la estrategia adoptada.

\section{Diagnóstico y entrega}
\begin{itemize}
\item Ante 401/403: comprobar sesión, proveedor, rol y políticas; no ampliar permisos
sin identificar la causa.
\item Si el registro vuelve a pedir confirmación: comprobar disponibilidad de
\code{registrar-cuenta} y si se activó el respaldo legacy.
\item Si falta correo: distinguir cuenta creada, solicitud al proveedor y recepción;
revisar registros sanitizados sin publicar datos o tokens.
\item Si una ruta directa falla: comprobar raíz, salida \code{www} y reescritura SPA.
\item Puerto ocupado: identificar el proceso propio antes de detenerlo; no cerrar
servidores ajenos.
\item Entrega: guardar versión, controles ejecutados, limitaciones, fuentes/PDF y
matrices. Adjuntar evidencia sanitizada y enlace de commit cuando el equipo lo autorice.
\end{itemize}
\end{document}
